% SPDX-License-Identifier: Apache-2.0
% covers: LinePair
\datasheet{LinePair}{The pixel-clock end of a scanout from memory: two
lines, one shown while the other fills, each asked for a line ahead.}

\dsfacts{\code{lib/parts/src/scanout.rs}}%
{\code{txhdl\_parts::scanout::LinePair<LEN, AW, ROWS, TOTAL, STRIDE, C>}}%
{\code{//docs:examples}}

\subsection*{Function}

A video peripheral that shows a frame from memory cannot wait for
memory pixel by pixel: the beam does not stop. \code{LinePair} sits on
the pixel clock beside the raster and holds two lines. The beam reads
one while the other fills, and they change places at the start of each
line. At that start it asks for the line after the one about to be
shown, so each line has the entire duration of the preceding line to
arrive. It asks for nothing until \code{show} is high, and then it asks for a
frame's first line, from the base, before any other. Out of a reset
the next line's address is zero, and the line there is the boot
memory's, whose one-beat answer to a burst hung the fetch for good on
the board (issue 1178). So the picture starts with the first frame
whose last row begins after \code{show} rises.
The bus side of the same scanout is \code{ScanFetch}, and between the
two the request and the words cross clocks through \code{ChanCdc}
(issue 151).

The frame's base address is a register taken at the vertical sync, so
a host that draws into one buffer and shows another flips them with
one write. A column the beam shows before its word has arrived sets
\code{starved}, a sticky bit that \code{clear} resets, so a run can
show that no line starved. The pair keeps the addresses of the lines
it has asked for and not had whole, at most four. A line that gets no
word for two line times in a row sets \code{stuck}, sticky too, and
puts its address on \code{stuck\_at}, so a fetch that hangs says so
rather than leaving a black screen and a clear \code{starved}
(issue 1197).

\subsection*{Parameters}

\begin{center}\footnotesize
\begin{tabular}{@{}lp{0.68\textwidth}@{}}
\toprule
Parameter & Meaning \\
\midrule
\code{LEN} & The words in a visible line, at most \code{1 << AW}. \\
\code{AW} & The width of a column. \\
\code{ROWS} & The visible rows of a frame. \\
\code{TOTAL} & The rows of a frame, blanking included. \\
\code{STRIDE} & The bytes from one line's start in memory to the next's. \\
\code{C} & The pixel clock. \\
\bottomrule
\end{tabular}
\end{center}

The tables below use \code{LinePair<8, 3, 4, 6, 32, ClkPix>}, the pair
of \code{ex\_scanpair}: lines of eight words, four visible rows of
six.

\dstables{LinePair}

\subsection*{Behaviour}

\begin{itemize}
\item A word is taken whenever one is offered, into the line filling
at the next column. The pair is what the fetch is backpressured by,
and a line longer than \code{LEN} loses its tail rather than stalling.
\item At a line's start (\code{line} high) the two lines change
places, and the line filling's count starts again. A word taken in
that same cycle counts in the line just filled.
\item At a line's start a request goes out on \code{req}, if it has
room: the address of the line after the one about to be shown. On the
frame's last row it is the frame's first line, at the base taken at
the last vertical sync.
\item While \code{show} is low nothing is asked for. The frame's first
line is asked for at the start of the frame's last row with
\code{show} high, and the lines after it only once it has been, so the
first request after a reset or a hide is always the base's line.
With \code{show} low the pair disarms: it asks for nothing more, and
does not set \code{starved}. Shown again, it starts from the base's
line as out of a reset.
\item \code{pix} is the word at \code{col} in the line not filling, a
cycle after \code{col} names it, or \code{LATE}, magenta
(\code{0x00ff00ff}), if that word has not arrived (issue 1209). No
picture is drawn in magenta, so a late column shows as itself rather
than as the word the line before left there.
\item A word belongs to the line it was asked for, which the lines owed
say (issue 1233). A line that starts before it has come whole is
\code{late}, and the rest of its words go into the half the beam
reads, so its columns fill in as they come. Lines owed whose time has
passed are counted in \code{behind}, and the rest of their words are
dropped. Neither lands in the half filling, where the next row would
show it as its own.
\item Once the frame's first line has been asked for, a visible column
at or past the words that line had when it was swapped in sets
\code{starved}, until \code{clear}.
\item A line asked for goes after the lines still owed, and no line is
asked for while four are. A word counts against the oldest owed, and
that line's last word takes it off.
\item Two line starts in a row with lines owed and no word come set
\code{stuck}, and \code{stuck\_at} is the oldest owed line's address
then; while \code{stuck} is set it does not change. \code{clear}
clears \code{stuck} and leaves \code{stuck\_at}. The pair goes on
asking.
\end{itemize}

\subsection*{Verification}

\begin{itemize}
\item \code{ex\_scanpair} runs \code{LinePair} on a pixel clock
unrelated to the bus clock, with \code{ScanFetch}, \code{LineFetch},
\code{AxiHost}, \code{AxiPer}, a memory and \code{ChanCdc} both ways.
The memory stalls its beats, standing in for a refresh. Three frames
with a short stall in every line show every pixel as the word at its
place and never set \code{starved}; a fourth with one stall longer than
a line sets it, and it stays set until it is cleared.
\item \code{//docs:scanpair\_sim\_linepair\_linepair\_tb\_test} and
\code{//docs:scanpair\_vsim\_linepair\_test} simulate the netlist
against that run under nvc and Verilator.
\end{itemize}

\subsection*{Used by}

The example \code{ex\_scanpair}, and \code{ScanVideo}, which the
flagship holds in its video slot (issue 151).

\subsection*{Limits}

One line a request, and a line is \code{LEN} words; there is no
horizontal offset or scaling. A column of a line that arrives late
shows \code{LATE} until its word comes, and \code{starved} says that
it happened; a line that never comes sets \code{stuck}. It runs on one clock, \code{C}.
